\begin{lemma}[The character at a totally split rational prime]\label{split:ray}
Let \(\mathfrak w\) lie above a rational prime completely split in the
Galois CM field \(E\). If its restriction to \(F\) is outside the
ramification set \(S\) of the triple, then
\(t=\nu_{\mathfrak w}(\pi)\) has infinite order, for either choice of
the place above \(F\) and every uniformizer \(\pi\).
\end{lemma}
\begin{proof}
Fix \(E\subset\mathbb C\). The chosen embeddings above the real places
form a CM type \(\Phi\), and we can regard them as automorphisms of
\(E\) by normality. Put \(m_\tau=m_\alpha(\tau)-m_\beta(\tau)>d\).
The infinite component of \(\nu\) at \(x\in E^\times\) is
\(\prod_{\tau\in\Phi}(\tau x/|\tau x|)^{m_\tau}\).
Take a finite ray modulus for \(\nu\), prime to \(\mathfrak w\).
By finiteness of the ray class group there are \(h>0\) and \(x\in
E^\times\) such that \((x)=\mathfrak w^h\) and \(x\equiv1\) modulo
this modulus. All ramified finite components are one at \(x\);
at every other finite place except \(\mathfrak w\), \(x\) is a unit.
The continuous Hecke character is trivial on diagonal \(E^\times\).
Its value on \(x\), squared, therefore gives
\[
 t^{-2h}=\prod_{\tau\in\Phi}
            \bigl(\tau x/\overline{\tau x}\bigr)^{m_\tau}.
\]
The principal fractional ideal on the right has exponents \(hm_\tau\)
and \(-hm_\tau\) at \(\tau\mathfrak w\) and
\(\bar\tau\mathfrak w\), respectively. These primes are all distinct
because the decomposition group is trivial. The fractional ideal is
nontrivial, whereas every root of unity generates the unit ideal.
This proves the assertion. Unramifiedness gives independence of \(\pi\).
\end{proof}
\src{src/Development/AbelianVarietyTheory/Definitions/Foundations/NotIsSplitOfIsInertAt.lean}{266}{Continuity and triviality on the diagonal};
\src{src/Development/AbelianVarietyTheory/BlockTheory/DivisorBase/WeilTensorModels.lean}{114}{Infinite character on ideles with trivial finite part};
\src{src/Development/AbelianVarietyTheory/BlockTheory/DivisorBase/WeilTensorModels.lean}{374}{Local ramification and infinite exponents}.

\begin{proposition}[The larger orbital transform at a good split place]
\label{split:good}
Suppose \(\nu\) is unramified and \(t=\nu(\pi)\) has infinite order.
Set \(n=a\), \(s=d/2\). On every matrix test \(\Phi\), the larger
transform in \eqref{newtheta:big}, initially integrated with exponent
\(z\) for \(\Re z>n\), has a regular continuation at \(z=s\).
This transform and the smaller transform \eqref{newtheta:small}
have the same image on their whole test spaces, namely the induced
representation in Proposition~\ref{newtheta:split}. The kernel on
each side is the span of its source-action differences. For
unramified \(\alpha,\beta\) their integral-matrix indicators can be
normalized to have the same nonzero image.
\end{proposition}
\begin{proof}
Fix a full-rank \(X\). The function \(H(g)=\Phi(gX)\) on \(M_n(k)\)
is locally constant and compactly supported, since multiplying by a
right inverse of \(X\) bounds its support. In additive matrix
coordinates the integral is a constant times
\[
 Z_n(H,z)=\int_{M_n(k)} H(g)\nu(\det g)|\det g|^{z-n}\,dg .
\]
Assign the integrand value zero on singular matrices, a null set
by successive one-variable polynomial testing and Fubini.
For \(0\le j\le n\) use the same exponent \(z-n\) in \(Z_j\) on
\(M_j(k)\), with \(Z_0\) evaluation at its single point.
We show by induction that \(Z_j\) is rational in \(q^{-z}\) and
its only possible poles for finite \(z\) satisfy
\[
           1-t^i q^{-i(z-n+i)}=0,\qquad 1\le i\le j .
\]
Away from a neighborhood of the zero matrix, a compact support has
a finite disjoint compact open subdivision on each part of which a
chosen nonzero pivot \(b\) has constant valuation. Permute rows and
columns to make it the top left entry. For the other three blocks use
coordinates \((B,C,D-Cb^{-1}B)\) instead of \((B,C,D)\). This
preserves additive measure by translation in \(D\) and factors the
determinant as \(b\det(D-Cb^{-1}B)\), up to a unit sign.
In these coordinates a compactly supported locally constant test is
a finite sum of products on compact open rectangles, by a finite
refinement of a cover by rectangles of constancy. The pivot weight
is a Laurent monomial times a constant. The remaining integrals thus
have the asserted form by induction.

Near zero subtract the constant value of the test times the
indicator of a sufficiently small origin ball. Scalar dilation
reduces that ball to \(M_j(\mathcal O)\) up to a Laurent monomial.
In this unit ball the complement of \(\pi M_j(\mathcal O)\) is
covered by the preceding pivot argument. The integral on the
inner ball is \(t^j q^{-j^2-j(z-n)}\) times the unit-ball integral.
Solving this identity in the half-plane of convergence introduces
exactly the indicated possible denominator with \(i=j\). This
proves rational continuation. For real \(z=s\), a zero of a
denominator would force first \(s-n+i=0\) by \(|t|=1\), and then
\(t^i=1\). Hence every test is regular there.

Source equivariance holds first in the convergent half-plane by
Haar substitution and hence at \(s\) by continuation. The same
reasoning gives the displayed parabolic covariance of
Proposition~\ref{newtheta:split} at \(s\); smoothness follows from
the compact open stabilizer of the test. On full-rank tests the
integrals are entire and the image and kernel were computed above.
On the rank-\(r<n\) stratum the stabilizer test in
Proposition~\ref{split:kernel} would require
\[
          |\det t_0|^r=\nu(\det t_0)|\det t_0|^s,\qquad
          t_0=\operatorname{diag}(I_r,\pi I_{n-r}).
\]
The moduli differ for \(r\ne s\), and for \(r=s\) equality would
give \(t^{n-r}=1\). Thus every singular stratum has zero quotient by
source differences. The rank filtration and the full-rank kernel
calculation prove the asserted kernel on all tests. The smaller
side is Proposition~\ref{split:kernel}.

In the convergent region the Hermite count in
Proposition~\ref{split:kernel} gives the larger spherical value
\(\prod_{j=0}^{n-1}(1-tq^{j-z})^{-1}\). Its continuation at \(s\)
is finite and nonzero; the smaller value has the same formula
with \(n=a'\). If \(\alpha,\beta\) are unramified both output
functions are right \(K_d\)-invariant. Such an induced function is
determined by its value at one since \(G_d=PK_d\), so dividing each
map by its indicated value identifies the two spherical outputs.
\end{proof}

The covariance is \(\delta_P^{1/2}
((\alpha\circ\det)\boxtimes(\beta\circ\det))\):
\(\delta_P(A,D')=|\det A|^{a'}|\det D'|^{-a}\), from conjugation on
the upper right block. The two determinant characters give an
irreducible unitary representation of the Levi, since their external
product is one-dimensional and unitary. Bernstein's irreducibility
theorem for normalized induction of irreducible unitary representations
of a Levi in \(\mathrm{GL}_d(k)\) therefore proves that this common
image is irreducible, both here and in Corollary~\ref{split:ram}.
Bernstein,
\emph{P-invariant distributions on \(\mathrm{GL}(N)\) and the classification
of unitary representations of \(\mathrm{GL}(N)\) (non-Archimedean case)},
\S0.2, corollary, proved in \S8.2~\cite{bernsteinInvariant}.

Lemma~\ref{split:ray} supplies the character hypothesis for every
good totally split place of the triple, at either splitting.
The construction here evaluates a rational continuation on unramified
tests. For the integral-matrix indicator infinitely many Cartan shells
have nonzero integral, so the finite-support sum defining the prescribed
shell map cannot replace this continuation. The matching condition
that specifies that shell map is imposed at \(S\), where the earlier
ramified calculation applies. The results here concern the displayed
matrix models. Passing to the given adelic Weil actions requires the
identifications with these polarizations and their lattice vacua.
\src{src/Development/AbelianVarietyTheory/Definitions/Divisors/RankEqCardIffExistsFrameMatrixFirstColumnBlock.lean}{707}{Finite-support Cartan sum in the prescribed map};
\src{src/Development/AbelianVarietyTheory/BlockTheory/DivisorBase/WeilTensorModels.lean}{470}{Places at which equality with that map is required}.
